% Part: conditional-logics
% Chapter: minimal-change-semantics
% Section: sphere-models

\documentclass[../../../include/open-logic-section]{subfiles}

\begin{document}

\olfileid[hi]{con}{min}{sph}

\olsection{गोलक प्रतिरूप}

प्रतितथ्यात्मक का औपचारिक अर्थविज्ञान देने का एक ढंग लुईस के अनौपचारिक विवरण
को गणितीय संरचना में बदलना है। तब किसी जगत~$w$ के चारों ओर के गोलक, जगतों
के समुच्चय हैं। चूँकि गोलक अंतर्निविष्ट हैं,~$w$ के चारों ओर जगतों के
समुच्चयों को उपसमुच्चय संबंध द्वारा रैखिक रूप से क्रमित होना चाहिए।

\begin{defn}
  कोई \emph{गोलक प्रतिरूप} त्रिक~$\mModel{M} = \tuple{W, O, V}$ है, जहाँ
  $W$ जगतों का अरिक्त समुच्चय, $V\colon \PVar \to \Pow{W}$ कोई मूल्यांकन,
  और $O\colon W \to \Pow{\Pow{W}}$ ऐसा फलन है जो प्रत्येक जगत~$w$ को
  \emph{गोलक-तंत्र}~$O_w$ देता है। प्रत्येक $w$ के लिए $O_w$, जगतों के
  समुच्चयों का समुच्चय है और उसे निम्न शर्तें संतुष्ट करनी चाहिए:
  \begin{enumerate}
  \item $O_w$,~$w$ पर \emph{केंद्रित} है: $\{w\} \in O_w$।
  \item $O_w$ \emph{अंतर्निविष्ट} है: जब भी $S_1$, $S_2 \in O_w$, तब
    $S_1 \subseteq S_2$ या $S_2 \subseteq S_1$; अर्थात $O_w$,~$\subseteq$
    द्वारा रैखिक रूप से क्रमित है।
  \item $O_w$ अरिक्त संघों के अधीन संवृत है।
  \item $O_w$ अरिक्त सर्वनिष्ठों के अधीन संवृत है।
  \end{enumerate}
\end{defn}

$O_w$ के पीछे की सहज धारणा यह है कि~$w$ के ``चारों ओर'' के जगतों को~$w$ से
उनकी दूरी के अनुसार स्तरों में रखा गया है। सबसे भीतरी गोलक में केवल स्वयं
$w$, अर्थात समुच्चय~$\{w\}$, है: किसी अन्य गोलक के जगतों की अपेक्षा $w$
स्वयं~$w$ के अधिक निकट है। यदि $S \subsetneq S'$, तो $S' \setminus S$ के
जगत~$S$ के जगतों की अपेक्षा~$w$ से अधिक दूर हैं: $S' \setminus S$,~$S$ और
~$S'$ के बाहर के जगतों के बीच की ``परत'' है। विशेषतः गोलकों को उनकी बाहरी
सतह के भीतर के सभी जगतों को समाविष्ट करने वाला समझना चाहिए; वे केवल अलग-अलग
परतें नहीं हैं।

\begin{figure}
\begin{center}
\begin{tikzpicture}[layerwidth=1.5,scale=.6]\tiny
  \spheresystem{3}
  \propositionintersect{0}{3}{60}{6}
  \path (0,0) node[world] {$w$};
  \spherepos{-30}{2}{node[world] {$w_2$}}
  \spherepos{220}{2}{node[world] {$w_3$}}
  \spherepos{90}{2}{node[world] {$w_1$}}
  \spherepos[shift={(.1,0)}]{10}{3}{node[world] {$w_5$}}
  \spherepos[shift={(.1,0)}]{-10}{3}{node[world] {$w_6$}}
  \spherepos{225}{3}{node[world] {$w_4$}}
  \spherepos{20}{4}{node[world] {$w_7$}}
  \path (5,0) node[left] {$p$};
\end{tikzpicture}
\caption{गोलक प्रतिरूप का आरेख}
\ollabel{fig:sphere-model}
\end{center}
\end{figure}

\olref{fig:sphere-model} का आरेख उस गोलक प्रतिरूप के अनुरूप है जिसमें
$W = \{w, w_1, \dots, w_7\}$ और $V(p) = \{w_5, w_6, w_7\}$। सबसे भीतरी
गोलक $S_1 = \{w\}$ है।~$w$ के निकटतम जगत $w_1$, $w_2$, $w_3$ हैं, अतः
अगला बड़ा गोलक $S_2 = \{w, w_1, w_2, w_3\}$ है। अधिक बाहर के जगत $w_4$,
$w_5$, $w_6$ हैं, अतः सबसे बाहरी गोलक
$S_3 = \{w, w_1, \dots, w_6\}$ है।~$w$ के चारों ओर गोलक-तंत्र
$O_w = \{S_1, S_2, S_3\}$ है। जगत~$w_7$,~$w$ के चारों ओर किसी गोलक में
नहीं है। वे निकटतम जगत जिनमें~$p$ सत्य है $w_5$ और~$w_6$ हैं, अतः सबसे
छोटा $p$-स्वीकारी गोलक~$S_3$ है।

किसी गोलक प्रतिरूप~$\mModel M$ में जगत~$w$ पर सूत्र~$!A$ की संतुष्टि,
$\mSat{M}{!A}[w]$, परिभाषित करने के लिए हम मोडल !!{formula}ों की परिभाषा में
$!B \cif !C$ का उपवाक्य जोड़ते हैं:

\begin{defn}
  $\mSat{M}{!B \cif !C}[w]$ तभी जब निम्न में से कोई एक हो:
  \begin{enumerate}
  \item\ollabel{sphere-vac} प्रत्येक $u \in \bigcup O_w$ के लिए
    $\mSat/{M}{!B}[u]$; अथवा
  \item\ollabel{sphere-nonvac} किसी $S \in O_w$ के लिए,
    \begin{enumerate}
      \item किसी $u \in S$ के लिए $\mSat{M}{!B}[u]$; और
      \item प्रत्येक $v \in S$ के लिए या तो
        $\mSat/{M}{!B}[v]$ या $\mSat{M}{!C}[v]$।
    \end{enumerate}
  \end{enumerate}
\end{defn}

इस परिभाषा के अनुसार $\mSat{M}{!B \cif !C}[w]$ तभी जब या तो पूर्ववर्ती~$!B$,
~$w$ के चारों ओर के गोलकों में सर्वत्र असत्य है, अथवा कोई गोलक~$S$ ऐसा है
जहाँ $!B$ सत्य और उस ``$!B$-स्वीकारी'' गोलक के सभी जगतों पर भौतिक सशर्त
$!B \lif !C$ सत्य है। ध्यान दें कि सहज व्याख्या से अपेक्षित बात के विपरीत,
हमने परिभाषा में यह अपेक्षा नहीं की कि $S$ \emph{सबसे भीतरी} $!B$-स्वीकारी
गोलक हो। किंतु यदि \olref{sphere-nonvac} की शर्त किसी गोलक~$S$ के लिए
संतुष्ट है, तो वह उन सभी गोलकों के लिए भी संतुष्ट है जिन्हें~$S$ समाविष्ट
करता है; अतः विशेषतः सबसे भीतरी गोलक के लिए।

यह भी ध्यान दें कि गोलक प्रतिरूपों की परिभाषा यह अपेक्षा नहीं करती कि कोई
सबसे भीतरी $!B$-स्वीकारी गोलक \emph{हो}: हमारे पास $!B$-स्वीकारी गोलकों का
अनंत अनुक्रम $S_1 \supsetneq S_2 \supsetneq \dots \supsetneq \{w\}$ हो
सकता है, और इसलिए कोई सबसे भीतरी $!B$-स्वीकारी गोलक नहीं। उस मामले में
$\mSat{M}{!B \cif !C}[w]$ तभी जब किसी~$i$ के लिए $!B \lif !C$, गोलकों
$S_i$, $S_{i+1}$, \dots में सर्वत्र सत्य हो।

\end{document}
